\section{Preliminaries}\label{sec:prelim}

\subsection{Machines and the statement (T)}\label{sec:prelim-machines}

We use the machine model of the first paper~\cite{FirstPaper2026}: a deterministic Turing machine with a finite state
set, a read-only two-way input tape, $k\ge1$ two-way work tapes over a finite alphabet, and a write-only one-way output
tape that receives at most one bit per step (the first paper also allows $k=0$; an idle work tape can always be added). All heads start at cell~$0$; head moves are $-1$, $0$ or $+1$; every step
writes the cell under each work head. For $d\ge1$ let $b(d)=\max(1,\clog d)$. The input is a $d\times d$ matrix with
$b(d)$-bit entries written row by row, of length $L=d^2b(d)$, and the machine must write its transpose in the same
encoding. Statement~(T) says that some machine does this within $c(L+1)$ steps for a constant~$c$ and all~$d$. Its
formal version, \lean{Transpose.LinearTimeTransposable}, is fixed in a frozen file of the first paper's library.

A machine is \emph{oblivious} if, for each input length, the positions of all heads at every time, the times at which it
appends output bits, and its halting time do not depend on the input. An oblivious machine has a single head schedule
per input length; an adaptive machine has one per input.

The Lean library of this paper does not model machines. Its theorems are about polymatroids, storage codes, matrices, and
labels on finite sets with hop bijections; the machine-level statements of \cref{sec:decomposition,sec:motion} are not formalized, and each comes
with a complete proof.

\subsection{Walks, hops and legs}\label{sec:prelim-walks}

A \emph{walk} of \emph{length}~$T$ is a map $x\colon\{0,\dots,T-1\}\to\Z$ with $\abs{x(t+1)-x(t)}\le1$. The work heads of a machine trace
walks $x_1,\dots,x_k$, one per work tape, and in a run these walks may be coupled arbitrarily.

For a head $h$ and a time $u$, let $\prev_h(u)=\max\{t<u : x_h(t)=x_h(u)\}$ when this set is nonempty. The pair
$(\prev_h(u),u)$ is the \emph{hop} of head $h$ ending at $u$: the symbol that head $h$ reads at step $u$ was written at
step $\prev_h(u)$, because every step writes the cell under each work head. A hop is \emph{long} if
$u-\prev_h(u)\ge2$; the other hops are pauses. The \emph{legs} of a walk partition its time interval greedily: the first
leg starts at time~$0$, each leg is the longest interval beginning at its first time on which the walk is weakly
monotone, and the next leg starts at the following time. So a leg ends just before the first step that reverses its direction.
A walk is \emph{home-based}, with home cell~$0$,
if every leg begins or ends within distance~$1$ of~$0$. Theorems about head motion are invariant under shifting a
tape, so the home cell can be any cell.

\subsection{Storage codes and guessing numbers}\label{sec:prelim-codes}

\emph{Side information} on a finite vertex set $V$ is a map $N\colon V\to 2^V$. For a simple graph $G$ on $V$ it is the
neighbourhood map $N(v)=\{u:uv\in E(G)\}$ (\lean{nbhd}); for a digraph it is the out-neighbourhood map
(\lean{outNbhd}). Let $\alpha$ be a finite alphabet.

\begin{definition}[Storage codes]\label{def:storage-code}
\formal{StorageCode.lean}{IsStorageCode, IsRecoverable}
A set $C\subseteq\alpha^V$ is a \emph{storage code} for $N$ if for all $x,y\in C$ and every $v\in V$, $x$ and $y$ agree at
$v$ whenever they agree on $N(v)$. A random word $X\colon\Omega\to\alpha^V$ is \emph{recoverable} for $N$ if the same holds
for all pairs of its values.
\end{definition}

So in a storage code each symbol is a function of the symbols at its neighbours. The \emph{rate} of $C$ is
$\log\card C/(\card V\log\card\alpha)$. The \emph{guessing number} $\gn(G,q)$ is the largest value of
$\log_q\card C$ over storage codes on $G$ with $\card\alpha=q$, and the asymptotic guessing number is
$\gn(G)=\sup_q\gn(G,q)$~\cite[Def.~3.7 and Thm.~3.6]{ChristofidesMarkstrom2011}; guessing numbers for a fixed alphabet go
back to Riis~\cite{Riis2007}. Storage codes are dual to index codes: for the
same side information, the optimal rates of the two problems are complementary asymptotically, and exactly for linear
codes~\cite{Mazumdar2014,ArbabjolfaeiKim2015}.
We work with storage codes throughout.
A \emph{linear} storage code over a field $F$ is a subspace of $F^V$ with the defining property; equivalently
(\cref{lem:linear-pattern}), a subspace of the kernel of a matrix $M\in F^{V\times V}$ with unit diagonal whose off-diagonal
entries vanish outside the edges of~$G$. The largest dimension of such a kernel is $\card V$ minus the
\emph{minrank} of $G$ over $F$~\cite{Haemers1979,BBJK2011}.

\subsection{Polymatroids and rate bounds}\label{sec:prelim-polymatroids}

\begin{definition}[Polymatroids, validity, rate bounds]\label{def:polymatroid}
\formal{Polymatroid.lean}{IsPolymatroid, IsValid, RateBound, IsGood}
A \emph{polymatroid} on $V$ is a function $h\colon 2^V\to\R$ with $h(\emptyset)=0$ that is monotone and submodular:
$h(A\cup B)+h(A\cap B)\le h(A)+h(B)$. It is \emph{valid} for side information $N$ if $h(N(v)\cup\{v\})=h(N(v))$ for every
$v$. Side information $N$ satisfies the \emph{rate bound} $c$ if every valid polymatroid satisfies
\[
  h(V)\;\le\;c\sum_{v\in V}h(\{v\}).
\]
A graph $G$ is \emph{good} if its neighbourhood map satisfies the rate bound~$\tfrac12$, that is, $2h(V)\le\sum_vh(\{v\})$ for
every valid polymatroid.
\end{definition}

The rate bound quantifies over all polymatroids, so it is a statement about the \emph{Shannon} information
inequalities; the smallest admissible $c$ is the Shannon bound $\lambda_{\Sh}$ of the side information.
For a random word $X$ on a finite probability space, $S\mapsto H(X_S)$ is a polymatroid (\cref{lem:entropic}),
and it is valid whenever $X$ is recoverable (the converse holds when every outcome has positive probability). Hence a rate bound $c$ bounds every storage code, linear or not, over
every alphabet: $H(X_V)\le c\sum_vH(X_v)$ for every recoverable $X$, and $\log\card C\le c\,\card V\log\card\alpha$
for every storage code $C$ (\cref{thm:entropic-rate}). Entropy is the finite Shannon entropy of the second
paper~\cite{SecondPaper2026}, whose Lean definitions we import.

\subsection{Formal status of the results}\label{sec:prelim-lean}

Every numbered result names its formal status in one of four ways. A line
\begin{center}
{\small\textsf{Lean:} \texttt{declaration, \dots}\hfill[\texttt{file}]}
\end{center}
gives the Lean declarations of a fully formalized result. The declarations live in the namespace
\texttt{MotionDepthCodes} and the files in the directory \leanfile{lean/MotionDepthCodes/}. The text after such a statement, labelled \emph{Outline of
the formal proof}, is a guide to the formal proof; it is a convenience, not a proof. A line beginning \textsf{Lean
(partial)} marks a result whose formal counterpart covers all but an explicitly named part, which the text proves or
cites. The words \textsf{Not formalized} mark a result with no formal counterpart; it comes with a complete proof,
labelled \emph{Proof (not formalized)}. A line beginning \textsf{Cited} marks a result of an earlier paper that is used
without proof, and says where it is proved and whether it is Lean-verified there. Where the English statement
simplifies, the Lean statement is authoritative.
Results from the literature are cited; a few short ones are reproved in our notation, and say so. Numerical
observations are labelled as such, claims that have only a proof sketch are stated in remarks beginning \emph{Proof sketch
only}, and neither is used in proofs. Numbered definitions carry a status line only when they
are formalized.

The library depends on Lean~4 and Mathlib (both v4.34.1) and on the Lean libraries of the first two papers of this
series, imported through relative paths. It contains no \texttt{sorry}, no \texttt{native\_decide} and no unsafe
evaluation, and every declaration cited here depends only on the axioms \texttt{propext}, \texttt{Classical.choice} and
\texttt{Quot.sound} (\cref{app:verify}). \Cref{app:index} lists every declaration of this library cited in the paper.
